\documentclass[11pt,a4paper]{article}

\usepackage[utf8]{inputenc}
\usepackage[T1]{fontenc}
\usepackage{amsmath,amssymb}
\usepackage{graphicx}
\usepackage{booktabs}
\usepackage{hyperref}
\usepackage[margin=0.9in]{geometry}
\usepackage{caption}
\usepackage{subcaption}
\usepackage{float}
\usepackage[dvipsnames]{xcolor}

\graphicspath{{../plotting/figures/}}

\title{Donut iso variants: inner-$R$ truth-purity scan and ABCD closure\\
\large All-calo 120\,MeV donut with $R_\mathrm{outer}=0.4$ (cross-check)}
\author{PPG12 Analysis Team (automated)}
\date{April 19, 2026}

\begin{document}
\maketitle

\section*{Executive summary}

Two new reco-isolation definitions were added as cross-check variants
(\texttt{syst\_type=None}, \texttt{syst\_role=None}):
\begin{itemize}
  \item \textbf{donutFull} (\texttt{use\_topo\_iso=4}):
    \texttt{cluster\_iso\_04} $-$ \texttt{cluster\_iso\_\{005,0075,02\}}.
    All-calo 120\,MeV tower sum (EMCal+HCalIn+HCalOut),
    outer $R=0.4$, inner $R \in \{0.05, 0.075, 0.20\}$.
    Cluster ET is subtracted as a scalar in both terms and cancels in the difference.
  \item \textbf{donutExcl} (\texttt{use\_topo\_iso=5}):
    \texttt{cluster\_iso\_excl\_04} $-$ \texttt{cluster\_iso\_excl\_\{005,0075,01,02\}}.
    Same three-calo 120\,MeV sum, but EMCal towers owned by the cluster are skipped at tower granularity
    rather than subtracted as a scalar. Inner $R \in \{0.05, 0.075, 0.10, 0.20\}$.
\end{itemize}

Switch added in \texttt{RecoEffCalculator\_TTreeReader.C:2042--2065}, mirrored in five companion macros.
\texttt{donutFull\_01} is deliberately skipped because \texttt{cluster\_iso\_01} does not exist in the slimtree
(logged in \texttt{reports/anatreemaker\_reprocess\_TODO.md}, along with the EMCal-only inner-ring branches
needed for a pure-EMCal donut).

\textbf{Key finding:} \texttt{donutFull\_02} and \texttt{donutExcl\_02} give
$\sim 2.5\times$ better ABCD-on-MC closure than nominal topo04 iso
(mean closure deviation $\langle|1-r|\rangle \approx 0.07$ vs.\ nominal $0.185$).
Low-pT closure is worst across \emph{all} variants (ratio $1.35$--$1.80$ at 9\,GeV),
high-pT closure is $0.83$--$0.87$ at 34\,GeV.

\section*{Two-pass workflow}

Per-variant iso cut was derived at 80\,\% target efficiency:
\begin{enumerate}
  \item \textbf{Pass 1}: run the full
    \texttt{RecoEffCalculator\_TTreeReader + MergeSim + CalculatePhotonYield}
    with placeholder loose cut ($b=100, s=0$).
  \item \textbf{Cut derivation}: parameterized
    \texttt{FindETCut.C(var\_type, target\_eff)} fits
    $E_T^\mathrm{iso,cut} = b + s\, p_T$ at 80\,\%;
    \texttt{derive\_iso\_cuts.py} drives all 7 variants.
  \item \textbf{Pass 2}: regenerate configs with derived $(b, s)$, re-run pipeline on HTCondor cluster 2701654.
\end{enumerate}

\begin{table}[H]
\centering
\caption{Derived 80\,\% iso cuts (isoET $< b + s\,p_T$).}
\begin{tabular}{lrr}
\toprule
Variant & $b$ [GeV] & $s$ \\
\midrule
donutFull\_005  & $0.1089$ & $0.0224$  \\
donutFull\_0075 & $0.2642$ & $0.0061$  \\
donutFull\_02   & $0.2305$ & $-0.0012$ \\
donutExcl\_005  & $0.2457$ & $0.0111$  \\
donutExcl\_0075 & $0.2774$ & $0.0053$  \\
donutExcl\_01   & $0.2921$ & $0.0017$  \\
donutExcl\_02   & $0.2305$ & $-0.0012$ \\
\bottomrule
\end{tabular}
\end{table}

\section*{Truth-matched signal purity}

Purity read from \texttt{Photon\_final\_bdt\_<var>\_mc.root} (MC-as-data, truth labels populated).

\begin{table}[H]
\centering
\caption{Truth-matched photon purity at three representative $p_T$ bins.
Nominal topo04 iso included for reference.}
\begin{tabular}{lrrr}
\toprule
Variant & $y(9\,\mathrm{GeV})$ & $y(21\,\mathrm{GeV})$ & $y(34\,\mathrm{GeV})$ \\
\midrule
nom (reference) & $0.581$ & $0.743$ & $0.797$ \\
donutFull\_005  & $0.684$ & $0.786$ & $0.849$ \\
donutFull\_0075 & $0.608$ & $0.805$ & $0.873$ \\
donutFull\_02   & $0.496$ & $0.728$ & $0.805$ \\
donutExcl\_005  & $0.648$ & $0.803$ & $0.858$ \\
donutExcl\_0075 & $0.606$ & $0.804$ & $0.874$ \\
donutExcl\_01   & $0.577$ & $0.793$ & $0.869$ \\
donutExcl\_02   & $0.496$ & $0.728$ & $0.805$ \\
\bottomrule
\end{tabular}
\end{table}

\begin{figure}[H]
\centering
\begin{subfigure}[t]{0.48\textwidth}
\includegraphics[width=\linewidth]{purity_sim_donutFull_scan.pdf}
\caption{donutFull family.}
\end{subfigure}\hfill
\begin{subfigure}[t]{0.48\textwidth}
\includegraphics[width=\linewidth]{purity_sim_donutExcl_scan.pdf}
\caption{donutExcl family.}
\end{subfigure}
\caption{Truth-matched signal purity vs.\ cluster $p_T$ for the donut-iso variants.
Tighter inner $R$ (smaller value) yields higher purity at low $p_T$; variants converge at high $p_T$.}
\end{figure}

\section*{isoET distribution by inner $R$ (representative $p_T$ bins)}

Overlay of the tight-shower-shape isoET distribution
(\texttt{h\_tight\_isoET\_0\_<ipt>}, from \texttt{data\_histo\_bdt\_<var>.root}),
normalized to unit area, across donut inner $R$ values at one low-$p_T$ bin
($10 < p_T < 12$\,GeV) and one higher-$p_T$ bin ($22 < p_T < 24$\,GeV).
Tight-region data is background-dominated below the iso cut; the distribution
shape reflects the underlying-event pedestal of the donut definition.

\begin{figure}[H]
\centering
\begin{subfigure}[t]{0.48\textwidth}
\includegraphics[width=\linewidth]{isoET_donutFull_ptlow.pdf}
\caption{donutFull, low $p_T$.}
\end{subfigure}\hfill
\begin{subfigure}[t]{0.48\textwidth}
\includegraphics[width=\linewidth]{isoET_donutFull_pthigh.pdf}
\caption{donutFull, high $p_T$.}
\end{subfigure}

\vspace{0.3em}

\begin{subfigure}[t]{0.48\textwidth}
\includegraphics[width=\linewidth]{isoET_donutExcl_ptlow.pdf}
\caption{donutExcl, low $p_T$.}
\end{subfigure}\hfill
\begin{subfigure}[t]{0.48\textwidth}
\includegraphics[width=\linewidth]{isoET_donutExcl_pthigh.pdf}
\caption{donutExcl, high $p_T$.}
\end{subfigure}
\caption{Tight-region isoET distributions in DATA for different donut inner $R$
at two representative $p_T$ bins. As $R_\mathrm{inner}$ grows, the annulus shrinks
and the distribution narrows / shifts toward zero. At low $p_T$ the distribution
is broader; at high $p_T$ it collapses toward the iso-cut edge.}
\end{figure}

\subsection*{Same in truth-matched MC signal (photon MC)}

\begin{figure}[H]
\centering
\begin{subfigure}[t]{0.48\textwidth}
\includegraphics[width=\linewidth]{isoET_donutFull_ptlow_sigMC.pdf}
\caption{donutFull, low $p_T$.}
\end{subfigure}\hfill
\begin{subfigure}[t]{0.48\textwidth}
\includegraphics[width=\linewidth]{isoET_donutFull_pthigh_sigMC.pdf}
\caption{donutFull, high $p_T$.}
\end{subfigure}

\vspace{0.3em}

\begin{subfigure}[t]{0.48\textwidth}
\includegraphics[width=\linewidth]{isoET_donutExcl_ptlow_sigMC.pdf}
\caption{donutExcl, low $p_T$.}
\end{subfigure}\hfill
\begin{subfigure}[t]{0.48\textwidth}
\includegraphics[width=\linewidth]{isoET_donutExcl_pthigh_sigMC.pdf}
\caption{donutExcl, high $p_T$.}
\end{subfigure}
\caption{Same as above for the truth-matched MC signal (merged photon MC,
\texttt{MC\_efficiency\_bdt\_<var>.root}). Distributions peak near zero and
are narrower than data, as expected for isolated prompt photons.
Shape differences across $R_\mathrm{inner}$ are subtler because the
underlying-event pedestal is smaller in signal-only MC.}
\end{figure}

\subsection*{Same in inclusive MC (photon + jet, weighted)}

\begin{figure}[H]
\centering
\begin{subfigure}[t]{0.48\textwidth}
\includegraphics[width=\linewidth]{isoET_donutFull_ptlow_incMC.pdf}
\caption{donutFull, low $p_T$.}
\end{subfigure}\hfill
\begin{subfigure}[t]{0.48\textwidth}
\includegraphics[width=\linewidth]{isoET_donutFull_pthigh_incMC.pdf}
\caption{donutFull, high $p_T$.}
\end{subfigure}

\vspace{0.3em}

\begin{subfigure}[t]{0.48\textwidth}
\includegraphics[width=\linewidth]{isoET_donutExcl_ptlow_incMC.pdf}
\caption{donutExcl, low $p_T$.}
\end{subfigure}\hfill
\begin{subfigure}[t]{0.48\textwidth}
\includegraphics[width=\linewidth]{isoET_donutExcl_pthigh_incMC.pdf}
\caption{donutExcl, high $p_T$.}
\end{subfigure}
\caption{Same as above for inclusive MC (photon MC $+$ jet MC, each weighted
by its cross section via MergeSim). This is the direct MC analogue of the
data plot above and should match data at cut1. The shape mirrors data closely,
confirming the donut definition's behavior on the background-dominated
tight-region population.}
\end{figure}

\section*{ABCD-on-MC closure}

The key validation: does the ABCD background subtraction reproduce the truth purity on MC?
The ratio $r = \mathrm{purity}_\mathrm{ABCD} / \mathrm{purity}_\mathrm{truth}$
is stored in \texttt{g\_mc\_purity\_fit\_ratio}; $r = 1$ is perfect closure.

\begin{table}[H]
\centering
\caption{ABCD/truth closure ratio at selected $p_T$ bins and mean absolute deviation from unity.
Lower $\langle|1-r|\rangle$ means better closure.}
\begin{tabular}{lrrrrrr}
\toprule
Variant & $r(9)$ & $r(15)$ & $r(21)$ & $r(28)$ & $r(34)$ & $\langle|1-r|\rangle$ \\
\midrule
nom (reference) & $1.795$ & $0.972$ & $0.910$ & $0.842$ & $0.832$ & $0.185$ \\
donutFull\_005  & $1.356$ & $0.972$ & $0.899$ & $0.826$ & $0.833$ & $0.136$ \\
donutFull\_0075 & $1.543$ & $1.004$ & $0.931$ & $0.861$ & $0.873$ & $0.138$ \\
\textbf{donutFull\_02} & $\mathbf{0.997}$ & $\mathbf{0.943}$ & $\mathbf{0.934}$ & $\mathbf{0.865}$ & $\mathbf{0.868}$ & $\mathbf{0.069}$ \\
donutExcl\_005  & $1.710$ & $1.004$ & $0.937$ & $0.844$ & $0.864$ & $0.164$ \\
donutExcl\_0075 & $1.536$ & $1.002$ & $0.929$ & $0.866$ & $0.874$ & $0.137$ \\
donutExcl\_01   & $1.346$ & $1.003$ & $0.922$ & $0.877$ & $0.872$ & $0.113$ \\
\textbf{donutExcl\_02} & $\mathbf{0.999}$ & $\mathbf{0.942}$ & $\mathbf{0.933}$ & $\mathbf{0.865}$ & $\mathbf{0.868}$ & $\mathbf{0.070}$ \\
\bottomrule
\end{tabular}
\end{table}

\begin{figure}[H]
\centering
\begin{subfigure}[t]{0.48\textwidth}
\includegraphics[width=\linewidth]{purity_closure_donutFull_scan.pdf}
\caption{donutFull family.}
\end{subfigure}\hfill
\begin{subfigure}[t]{0.48\textwidth}
\includegraphics[width=\linewidth]{purity_closure_donutExcl_scan.pdf}
\caption{donutExcl family.}
\end{subfigure}
\caption{Closure of ABCD-derived purity against truth-matched purity on MC.
The dashed line at 1.0 is perfect closure. All variants (including nominal) show
a common low-$p_T$ over-prediction and high-$p_T$ under-prediction;
inner $R = 0.20$ variants are closest to unity across the full $p_T$ range.}
\end{figure}

\section*{Physics observations}

\begin{itemize}
  \item \textbf{Closure is the headline result.} The nominal topo04 iso
    has an 18.5\% average deviation from unity; donut variants with
    $R_\mathrm{inner} = 0.20$ reduce this to 6.9--7.0\%. Other donut
    variants fall in between (11--16\%).
  \item \textbf{Low-$p_T$ non-closure is universal.}
    At $p_T = 9$\,GeV the ratio is $1.35$--$1.80$ across \emph{all}
    variants including nominal. This is probably driven by the fit-region
    statistics at the edge of the reco $p_T$ binning, not by the iso
    definition itself.
  \item \textbf{High-$p_T$ under-prediction is also universal}, with
    $r \approx 0.83$--$0.87$ at 34\,GeV across all variants. ABCD consistently
    under-estimates truth purity at high $p_T$ regardless of the iso
    definition.
  \item \textbf{Inner $R$ dominates purity, not outer-cone treatment.}
    donutFull vs.\ donutExcl at matched $R_\mathrm{inner}$ give truth
    purities within $\pm 0.05$ at low $p_T$ and $\lesssim 0.02$ at high
    $p_T$. \texttt{donutFull\_02} and \texttt{donutExcl\_02} are
    numerically identical to four decimals because the fitted iso cut is
    the same (see derived $(b, s)$).
  \item \textbf{Tighter inner $R$ gives higher truth purity at low $p_T$.}
    For the Excl family, $y(9\,\mathrm{GeV})$ monotonically decreases
    from $0.648$ at $R_\mathrm{inner} = 0.05$ to $0.496$ at $R_\mathrm{inner} = 0.20$.
    Interpretation: larger inner $R$ admits less annulus area, giving a
    looser absolute iso cut that keeps more background.
\end{itemize}

\section*{Caveats and follow-ups}

\begin{itemize}
  \item \textbf{Outer cone includes HCal.} The donut is not pure EMCal.
    A pure-EMCal donut requires new inner-ring branches
    \texttt{cluster\_iso\_\{005,0075,01,02\}\_emcal} at matching threshold
    (logged in \texttt{reports/anatreemaker\_reprocess\_TODO.md}).
  \item \textbf{\texttt{donutFull\_01} skipped} pending upstream
    \texttt{cluster\_iso\_01} branch (all-calo, 120\,MeV, $R<0.1$).
  \item \textbf{Degraded mirrors} (\texttt{Cluster\_rbr.C},
    \texttt{EtaMigrationStudy.C}, \texttt{compare\_energy\_response.C})
    silently use the full outer cone without inner subtraction for cases 4/5
    since they do not read \texttt{iso\_emcalinnerr}. Do not run donut
    variants through these diagnostic utilities expecting donut semantics.
  \item \textbf{MC iso scale/shift} set to 1.0/0.0 for all donut variants
    because the nominal 1.2/0.1 was tuned against full-cone topo04. Retuning
    would be needed if these ever became systematic variants.
  \item \textbf{The low-$p_T$ closure deviation is present in nominal too.}
    This is a pre-existing effect of the ABCD method at the binning edge,
    independent of the iso definition. Worth understanding before any
    purity systematic is claimed.
\end{itemize}

\section*{Reviewer pass matrix}

\begin{table}[H]
\centering
\begin{tabular}{lcccc}
\toprule
Artifact & Physics & Cosmetics & Re-derivation & Fix-validator \\
\midrule
\texttt{use\_topo\_iso=4,5} switch & PASS (2 WARN util) & N/A & N/A & PASS \\
donut truth-purity scan PDFs & N/A & PASS (2 WARN) & PASS & N/A \\
donut closure PDFs & N/A & (not yet reviewed) & PASS & N/A \\
80\,\% iso cuts & N/A & N/A & PASS & PASS \\
\texttt{g\_purity\_truth} (7 variants) & N/A & N/A & PASS & PASS \\
\texttt{g\_mc\_purity\_fit\_ratio} (7 variants) & N/A & N/A & PASS & PASS \\
\bottomrule
\end{tabular}
\end{table}



\section*{Per-variant MC purity closure}

Each panel below shows, for one variant, the ABCD-derived purity on MC (points with error bars),
the fit to the ABCD points (solid line), the signal-leakage-corrected ABCD purity,
and the truth-matched purity (red markers). For a closure-good variant, the ABCD
points/fit should track the truth markers; deviations mark where the ABCD
subtraction biases the purity estimate at that $p_T$.
Nominal topo04 iso (\texttt{bdt\_nom}) is included as a reference.

\begin{figure}[H]
\centering
\begin{subfigure}[t]{0.24\textwidth}
\includegraphics[width=\linewidth]{purity_sim_bdt_nom.pdf}
\caption{nom (reference).}
\end{subfigure}%
\begin{subfigure}[t]{0.24\textwidth}
\includegraphics[width=\linewidth]{purity_sim_bdt_donutFull_005.pdf}
\caption{donutFull\_005.}
\end{subfigure}%
\begin{subfigure}[t]{0.24\textwidth}
\includegraphics[width=\linewidth]{purity_sim_bdt_donutFull_0075.pdf}
\caption{donutFull\_0075.}
\end{subfigure}%
\begin{subfigure}[t]{0.24\textwidth}
\includegraphics[width=\linewidth]{purity_sim_bdt_donutFull_02.pdf}
\caption{donutFull\_02.}
\end{subfigure}

\vspace{0.3em}

\begin{subfigure}[t]{0.24\textwidth}
\includegraphics[width=\linewidth]{purity_sim_bdt_donutExcl_005.pdf}
\caption{donutExcl\_005.}
\end{subfigure}%
\begin{subfigure}[t]{0.24\textwidth}
\includegraphics[width=\linewidth]{purity_sim_bdt_donutExcl_0075.pdf}
\caption{donutExcl\_0075.}
\end{subfigure}%
\begin{subfigure}[t]{0.24\textwidth}
\includegraphics[width=\linewidth]{purity_sim_bdt_donutExcl_01.pdf}
\caption{donutExcl\_01.}
\end{subfigure}%
\begin{subfigure}[t]{0.24\textwidth}
\includegraphics[width=\linewidth]{purity_sim_bdt_donutExcl_02.pdf}
\caption{donutExcl\_02.}
\end{subfigure}
\caption{Per-variant MC purity closure plots. Each panel: ABCD-derived purity points (from
\texttt{gpurity}/\texttt{gpurity\_leak}), the ABCD fit, and truth-matched purity
(\texttt{g\_purity\_truth}). The quality of agreement between ABCD and truth across $p_T$
is the closure check summarized in the previous section's ratio table.
\texttt{donutFull\_02} and \texttt{donutExcl\_02} show the tightest truth--ABCD agreement,
especially at low $p_T$.}
\end{figure}



\section*{Final cross section per variant}

\subsection*{Donut / nominal ratio}

The following ratio is defined per reco $p_T$ bin on the unfolded signal yield
\texttt{h\_unfold\_sub\_result} (which is the background-subtracted, unfolded
yield used to compute $d^2\sigma/d\eta\, dp_T$): $r = \mathrm{yield}_\mathrm{donut} /
\mathrm{yield}_\mathrm{nom}$.

\begin{table}[H]
\centering
\caption{Donut / nominal unfolded-yield ratio at selected $p_T$ bins.
Mean absolute deviation from unity averaged over the 12 reco $p_T$ bins.}
\begin{tabular}{lrrrrrr}
\toprule
Variant & $r(9)$ & $r(15)$ & $r(21)$ & $r(28)$ & $r(34)$ & $\langle|1-r|\rangle$ \\
\midrule
donutFull\_005  & $1.129$ & $0.929$ & $0.844$ & $0.818$  & $0.788$  & $0.130$ \\
donutFull\_0075 & $1.035$ & $0.921$ & $0.837$ & $0.903$  & $0.835$  & $0.097$ \\
donutFull\_02   & $1.354$ & $0.974$ & $0.999$ & $1.310$  & $1.048$  & $0.138$ \\
\textbf{donutExcl\_005 (anomalous)} & $0.812$ & $0.908$ & $1.937$ & $\mathbf{-6.02}$ & $\mathbf{-0.04}$ & $1.665$ \\
donutExcl\_0075 & $1.042$ & $0.917$ & $0.830$ & $0.869$  & $0.837$  & $0.104$ \\
donutExcl\_01   & $1.137$ & $0.938$ & $0.882$ & $1.018$  & $0.822$  & $0.096$ \\
donutExcl\_02   & $1.352$ & $0.974$ & $0.999$ & $1.310$  & $1.048$  & $0.137$ \\
\bottomrule
\end{tabular}
\end{table}

\textbf{donutExcl\_005 anomaly.} This variant yields negative unfolded
signal at $p_T \ge 28$\,GeV, implying the ABCD background estimate exceeds
the tight+iso count in those bins. Looking at the raw
\texttt{h\_unfold\_sub\_result} across $p_T$: values at $p_T = 22$--$26$
are anomalously large (2--6$\times$ nominal) before the bin-11 sign flip to $-5.66$
and bin-12 $-0.81$. Truth purity for the same variant is healthy (0.65--0.86),
so the pathology is in yield extraction, not purity. Likely an interaction
between the tight-iso count and the R-factor estimator at low statistics
with the tightest inner $R$ + ownership-excluded outer cone. Recommend:
re-run with higher MC statistics (more condor jobs per sample) or drop this
variant from the scan.

\begin{figure}[H]
\centering
\begin{subfigure}[t]{0.48\textwidth}
\includegraphics[width=\linewidth]{xsec_ratio_donutFull_scan.pdf}
\caption{donutFull / nominal.}
\end{subfigure}\hfill
\begin{subfigure}[t]{0.48\textwidth}
\includegraphics[width=\linewidth]{xsec_ratio_donutExcl_scan.pdf}
\caption{donutExcl / nominal.}
\end{subfigure}
\caption{Unfolded-yield ratio (donut / nominal) vs.\ reco $p_T$.
Dashed line at unity. Excluding \texttt{donutExcl\_005}, variants agree with
nominal to $\sim 10$--$14\%$ on $\langle|1-r|\rangle$. The \texttt{R\_inner}=0.20
variants diverge upward at low $p_T$ (ratio up to 1.35) and stay within
$\pm 5$\% of nominal at mid-$p_T$. \texttt{donutExcl\_005} is plotted but
pT $\ge 26$\,GeV points fall off-scale (negative / $\gg 1.4$).}
\end{figure}

\subsection*{Per-variant final cross section}

\begin{figure}[H]
\centering
\begin{subfigure}[t]{0.24\textwidth}
\includegraphics[width=\linewidth]{final_bdt_nom.pdf}
\caption{nom.}
\end{subfigure}%
\begin{subfigure}[t]{0.24\textwidth}
\includegraphics[width=\linewidth]{final_bdt_donutFull_005.pdf}
\caption{donutFull\_005.}
\end{subfigure}%
\begin{subfigure}[t]{0.24\textwidth}
\includegraphics[width=\linewidth]{final_bdt_donutFull_0075.pdf}
\caption{donutFull\_0075.}
\end{subfigure}%
\begin{subfigure}[t]{0.24\textwidth}
\includegraphics[width=\linewidth]{final_bdt_donutFull_02.pdf}
\caption{donutFull\_02.}
\end{subfigure}

\vspace{0.3em}

\begin{subfigure}[t]{0.24\textwidth}
\includegraphics[width=\linewidth]{final_bdt_donutExcl_005.pdf}
\caption{donutExcl\_005 (anomalous).}
\end{subfigure}%
\begin{subfigure}[t]{0.24\textwidth}
\includegraphics[width=\linewidth]{final_bdt_donutExcl_0075.pdf}
\caption{donutExcl\_0075.}
\end{subfigure}%
\begin{subfigure}[t]{0.24\textwidth}
\includegraphics[width=\linewidth]{final_bdt_donutExcl_01.pdf}
\caption{donutExcl\_01.}
\end{subfigure}%
\begin{subfigure}[t]{0.24\textwidth}
\includegraphics[width=\linewidth]{final_bdt_donutExcl_02.pdf}
\caption{donutExcl\_02.}
\end{subfigure}
\caption{Per-variant final isolated-photon cross section
$d^2\sigma/d\eta\, dp_T$ vs.\ $p_T^\gamma$, produced by
\texttt{plot\_final\_selection.C}. Nominal topo04 iso (first panel) for reference.
The \texttt{donutExcl\_005} panel shows the high-$p_T$ anomaly noted above.}
\end{figure}



\section*{Isolation ET template fits (CONF)}

Produced by \texttt{plotting/CONF\_plots.C} from
\texttt{data\_histo\_bdt\_<var>.root} (tight and non-tight data) and
\texttt{MC\_efficiency\_bdt\_<var>.root} (MC signal template). Two $p_T$ bin
ranges shown as representatives: \textbf{8--14\,GeV} (bins 0--2) and
\textbf{20--26\,GeV} (bins 6--8). The fit overlays the MC signal shape + a
non-tight-derived background template onto the tight-region data,
extracting the signal fraction per pT range. Full 4-range sets
(0--2, 3--5, 6--8, 9--9) for all 7 donut variants are in
\texttt{plotting/figures/h1D\_iso\_fit\_bdt\_donut*.pdf}; the selection
report includes them in each per-config section via the standard PLOT\_GROUPS
``Isolation fit (CONF)'' entry.

\subsection*{Low $p_T$ bin range: $8 < p_T < 14$\,GeV (bins 0--2)}

\begin{figure}[H]
\centering
\begin{subfigure}[t]{0.24\textwidth}
\includegraphics[width=\linewidth]{h1D_iso_bdt_nom_0_2.pdf}
\caption{nom.}
\end{subfigure}%
\begin{subfigure}[t]{0.24\textwidth}
\includegraphics[width=\linewidth]{h1D_iso_bdt_donutFull_005_0_2.pdf}
\caption{donutFull\_005.}
\end{subfigure}%
\begin{subfigure}[t]{0.24\textwidth}
\includegraphics[width=\linewidth]{h1D_iso_bdt_donutFull_0075_0_2.pdf}
\caption{donutFull\_0075.}
\end{subfigure}%
\begin{subfigure}[t]{0.24\textwidth}
\includegraphics[width=\linewidth]{h1D_iso_bdt_donutFull_02_0_2.pdf}
\caption{donutFull\_02.}
\end{subfigure}

\vspace{0.3em}

\begin{subfigure}[t]{0.24\textwidth}
\includegraphics[width=\linewidth]{h1D_iso_bdt_donutExcl_005_0_2.pdf}
\caption{donutExcl\_005.}
\end{subfigure}%
\begin{subfigure}[t]{0.24\textwidth}
\includegraphics[width=\linewidth]{h1D_iso_bdt_donutExcl_0075_0_2.pdf}
\caption{donutExcl\_0075.}
\end{subfigure}%
\begin{subfigure}[t]{0.24\textwidth}
\includegraphics[width=\linewidth]{h1D_iso_bdt_donutExcl_01_0_2.pdf}
\caption{donutExcl\_01.}
\end{subfigure}%
\begin{subfigure}[t]{0.24\textwidth}
\includegraphics[width=\linewidth]{h1D_iso_bdt_donutExcl_02_0_2.pdf}
\caption{donutExcl\_02.}
\end{subfigure}
\caption{Isolation-ET templates vs.\ tight data (points), with
MC signal template (blue fill) and non-tight data-derived background template
(red fill), for $8 < p_T < 14$\,GeV. Nominal reference and all 7 donut variants.}
\end{figure}

\subsection*{High $p_T$ bin range: $20 < p_T < 26$\,GeV (bins 6--8)}

\begin{figure}[H]
\centering
\begin{subfigure}[t]{0.24\textwidth}
\includegraphics[width=\linewidth]{h1D_iso_bdt_nom_6_8.pdf}
\caption{nom.}
\end{subfigure}%
\begin{subfigure}[t]{0.24\textwidth}
\includegraphics[width=\linewidth]{h1D_iso_bdt_donutFull_005_6_8.pdf}
\caption{donutFull\_005.}
\end{subfigure}%
\begin{subfigure}[t]{0.24\textwidth}
\includegraphics[width=\linewidth]{h1D_iso_bdt_donutFull_0075_6_8.pdf}
\caption{donutFull\_0075.}
\end{subfigure}%
\begin{subfigure}[t]{0.24\textwidth}
\includegraphics[width=\linewidth]{h1D_iso_bdt_donutFull_02_6_8.pdf}
\caption{donutFull\_02.}
\end{subfigure}

\vspace{0.3em}

\begin{subfigure}[t]{0.24\textwidth}
\includegraphics[width=\linewidth]{h1D_iso_bdt_donutExcl_005_6_8.pdf}
\caption{donutExcl\_005.}
\end{subfigure}%
\begin{subfigure}[t]{0.24\textwidth}
\includegraphics[width=\linewidth]{h1D_iso_bdt_donutExcl_0075_6_8.pdf}
\caption{donutExcl\_0075.}
\end{subfigure}%
\begin{subfigure}[t]{0.24\textwidth}
\includegraphics[width=\linewidth]{h1D_iso_bdt_donutExcl_01_6_8.pdf}
\caption{donutExcl\_01.}
\end{subfigure}%
\begin{subfigure}[t]{0.24\textwidth}
\includegraphics[width=\linewidth]{h1D_iso_bdt_donutExcl_02_6_8.pdf}
\caption{donutExcl\_02.}
\end{subfigure}
\caption{Same as above for $20 < p_T < 26$\,GeV. The \texttt{donutExcl\_005}
high-$p_T$ pathology noted in the cross-section section shows up here as
the signal+background templates failing to cover the tight-region data shape.}
\end{figure}



\section*{BDT model alternatives: \texttt{base\_v0} and \texttt{base\_v0E}}

Cross-check on BDT model choice: two additional variants apply earlier model
versions (\texttt{base\_v0}, \texttt{base\_v0E}) across both ET bins via
\texttt{bdt\_et\_bin\_models=["base\_v0","base\_v0"]} (resp.\ v0E). Nominal
topo04 iso and nominal tight/non-tight BDT cuts are kept. Both $cluster\_bdt\_CLUSTERINFO\_CEMC\_base\_v0$ and $*\_base\_v0E$
branches already exist in the slimtree; no re-apply needed.
HTCondor cluster 2734512 (2 jobs, $\sim$20\,min each).

\subsection*{Summary numbers}

\begin{table}[H]
\centering
\caption{Truth purity, ABCD/truth closure, and unfolded-yield ratio to nominal,
at selected $p_T$ bins. Lower $\langle|1-r|\rangle$ is better closure.
Higher $|1-r_\mathrm{xsec}|$ means larger deviation of the cross section from nominal.}
\begin{tabular}{lrrrrrr}
\toprule
 & \multicolumn{3}{c}{Truth purity} & \multicolumn{3}{c}{ABCD/truth closure} \\
\cmidrule(lr){2-4}\cmidrule(lr){5-7}
Variant & $y(9)$ & $y(21)$ & $y(34)$ & $r(9)$ & $r(21)$ & $\langle|1-r|\rangle$ \\
\midrule
nom (reference)       & $0.581$ & $0.743$ & $0.797$ & $1.795$ & $0.910$ & $0.185$ \\
bdtmodel\_v0          & $0.544$ & $0.744$ & $0.782$ & $2.177$ & $0.973$ & $0.226$ \\
bdtmodel\_v0E         & $0.557$ & $0.753$ & $0.786$ & $2.491$ & $0.961$ & $0.265$ \\
\bottomrule
\end{tabular}
\end{table}

\begin{table}[H]
\centering
\caption{Unfolded-yield ratio (variant / nominal) across the $p_T$ spectrum.}
\begin{tabular}{lrrrrrr}
\toprule
Variant & $r(9)$ & $r(15)$ & $r(21)$ & $r(28)$ & $r(34)$ & $\langle|1-r|\rangle$ \\
\midrule
bdtmodel\_v0    & $0.751$ & $0.980$ & $1.014$ & $1.183$ & $1.604$ & $0.196$ \\
bdtmodel\_v0E   & $0.760$ & $0.974$ & $1.006$ & $1.219$ & $1.742$ & $0.223$ \\
\bottomrule
\end{tabular}
\end{table}

\textbf{Observations.}
(1) Truth purity for the v0/v0E models is close to (but slightly lower than)
nominal at low $p_T$ ($0.54$--$0.56$ vs.\ $0.58$), and essentially identical
at high $p_T$.
(2) \textbf{ABCD closure is worse} for both alternatives:
$\langle|1-r|\rangle = 0.226$ (v0) and $0.265$ (v0E) vs.\ $0.185$ nominal.
The low-$p_T$ over-prediction is particularly strong ($r = 2.18$ and $2.49$
at $p_T = 9$\,GeV vs.\ $1.80$ nominal), suggesting the ABCD R-factor
assumption is less satisfied with these BDT discriminants.
(3) Cross section at high $p_T$ diverges sharply from nominal:
$1.60\times$ and $1.74\times$ at $p_T=34$\,GeV respectively.
The v0/v0E models either admit substantially more tight+iso yield
(looser BDT) or have larger non-closure absorbing into the subtracted
``signal'' yield. Given that the cross-section ratio grows with $p_T$ while
closure degrades at low $p_T$, the high-$p_T$ excess is likely driven by
poorer signal--background separation of the older BDT variants in the
region where statistics are thin.

\subsection*{Per-variant MC purity closure}

\begin{figure}[H]
\centering
\begin{subfigure}[t]{0.33\textwidth}
\includegraphics[width=\linewidth]{purity_sim_bdt_nom.pdf}
\caption{nom (reference).}
\end{subfigure}%
\begin{subfigure}[t]{0.33\textwidth}
\includegraphics[width=\linewidth]{purity_sim_bdt_bdtmodel_v0.pdf}
\caption{bdtmodel\_v0.}
\end{subfigure}%
\begin{subfigure}[t]{0.33\textwidth}
\includegraphics[width=\linewidth]{purity_sim_bdt_bdtmodel_v0E.pdf}
\caption{bdtmodel\_v0E.}
\end{subfigure}
\caption{MC purity closure for the two alternative BDT models. Each panel
shows the ABCD-derived purity (points), the fit, and the truth-matched
purity (red markers).}
\end{figure}

\subsection*{Isolation ET template fits}

\begin{figure}[H]
\centering
\begin{subfigure}[t]{0.24\textwidth}
\includegraphics[width=\linewidth]{h1D_iso_bdt_nom_0_2.pdf}
\caption{nom, $8<p_T<14$.}
\end{subfigure}%
\begin{subfigure}[t]{0.24\textwidth}
\includegraphics[width=\linewidth]{h1D_iso_bdt_nom_6_8.pdf}
\caption{nom, $20<p_T<26$.}
\end{subfigure}%
\begin{subfigure}[t]{0.24\textwidth}
\includegraphics[width=\linewidth]{h1D_iso_bdt_bdtmodel_v0_0_2.pdf}
\caption{v0, $8<p_T<14$.}
\end{subfigure}%
\begin{subfigure}[t]{0.24\textwidth}
\includegraphics[width=\linewidth]{h1D_iso_bdt_bdtmodel_v0_6_8.pdf}
\caption{v0, $20<p_T<26$.}
\end{subfigure}

\vspace{0.3em}

\begin{subfigure}[t]{0.24\textwidth}
\includegraphics[width=\linewidth]{h1D_iso_bdt_bdtmodel_v0E_0_2.pdf}
\caption{v0E, $8<p_T<14$.}
\end{subfigure}%
\begin{subfigure}[t]{0.24\textwidth}
\includegraphics[width=\linewidth]{h1D_iso_bdt_bdtmodel_v0E_6_8.pdf}
\caption{v0E, $20<p_T<26$.}
\end{subfigure}%
\begin{subfigure}[t]{0.24\textwidth}
\includegraphics[width=\linewidth]{final_bdt_bdtmodel_v0.pdf}
\caption{Final $d^2\sigma$, v0.}
\end{subfigure}%
\begin{subfigure}[t]{0.24\textwidth}
\includegraphics[width=\linewidth]{final_bdt_bdtmodel_v0E.pdf}
\caption{Final $d^2\sigma$, v0E.}
\end{subfigure}
\caption{CONF template overlays (\texttt{h1D\_iso\_*}: tight data points,
MC signal template in blue fill, non-tight background template in red fill)
for nominal and the two alternative BDT models at two $p_T$ ranges,
followed by the unfolded final cross section per alternative model.
Poorer template agreement for v0/v0E at high $p_T$ is consistent with the
weaker closure reported above.}
\end{figure}



\section*{Final cross section with MC purity correction turned on}

Same 7 donut variants as Figure~8, now run with \texttt{mc\_purity\_correction=1}
in the config (the existing \texttt{config\_bdt\_mc\_purity\_correction.yaml}
mechanism applied per-variant). In this mode \texttt{CalculatePhotonYield.C}
uses the MC-as-data truth purity (from \texttt{Photon\_final\_bdt\_<var>\_mcpc\_mc.root})
as a multiplicative correction when extracting the data purity, rather than
trusting the ABCD fit alone. Truth purity and ABCD-on-MC closure are unchanged
from the non-mcpc run (they are MC-only quantities); the difference is in the
data yield.

\begin{table}[H]
\centering
\caption{mcpc: cross-section ratio (variant / nominal, both without mcpc for the
nominal denominator) and mean absolute deviation.
Compare with the corresponding non-mcpc row from the earlier cross-section table.}
\begin{tabular}{lrrrrr}
\toprule
Variant & $r_\mathrm{xs}(9)$ & $r_\mathrm{xs}(21)$ & $r_\mathrm{xs}(34)$ & $\langle|1-r_\mathrm{xs}|\rangle$ \\
\midrule
donutFull\_005\_mcpc  & $1.534$ & $0.844$ & $0.655$ & $0.282$ \\
donutFull\_0075\_mcpc & $1.599$ & $0.837$ & $0.725$ & $0.257$ \\
donutFull\_02\_mcpc   & $1.352$ & $0.999$ & $0.910$ & $0.129$ \\
donutExcl\_005\_mcpc  & $1.619$ & $0.937$ & $0.673$ & $0.321$ \\
donutExcl\_0075\_mcpc & $1.607$ & $0.929$ & $0.726$ & $0.264$ \\
donutExcl\_01\_mcpc   & $1.540$ & $0.922$ & $0.714$ & $0.220$ \\
donutExcl\_02\_mcpc   & $1.350$ & $0.999$ & $0.911$ & $0.129$ \\
\bottomrule
\end{tabular}
\end{table}

\textbf{Note on \texttt{donutExcl\_005}.} Without MC purity correction this
variant showed negative unfolded yields at $p_T \ge 28$\,GeV ($\langle|1-r|\rangle = 1.665$).
With mcpc on, the yields are all positive and $\langle|1-r|\rangle$ drops to
$0.321$. MC purity correction effectively rescues the pathological variant
by using the MC truth purity in place of the badly-behaved ABCD purity at high $p_T$.

\begin{figure}[H]
\centering
\begin{subfigure}[t]{0.24\textwidth}
\includegraphics[width=\linewidth]{final_bdt_nom.pdf}
\caption{nom (no mcpc, reference).}
\end{subfigure}%
\begin{subfigure}[t]{0.24\textwidth}
\includegraphics[width=\linewidth]{final_bdt_donutFull_005_mcpc.pdf}
\caption{donutFull\_005\_mcpc.}
\end{subfigure}%
\begin{subfigure}[t]{0.24\textwidth}
\includegraphics[width=\linewidth]{final_bdt_donutFull_0075_mcpc.pdf}
\caption{donutFull\_0075\_mcpc.}
\end{subfigure}%
\begin{subfigure}[t]{0.24\textwidth}
\includegraphics[width=\linewidth]{final_bdt_donutFull_02_mcpc.pdf}
\caption{donutFull\_02\_mcpc.}
\end{subfigure}

\vspace{0.3em}

\begin{subfigure}[t]{0.24\textwidth}
\includegraphics[width=\linewidth]{final_bdt_donutExcl_005_mcpc.pdf}
\caption{donutExcl\_005\_mcpc.}
\end{subfigure}%
\begin{subfigure}[t]{0.24\textwidth}
\includegraphics[width=\linewidth]{final_bdt_donutExcl_0075_mcpc.pdf}
\caption{donutExcl\_0075\_mcpc.}
\end{subfigure}%
\begin{subfigure}[t]{0.24\textwidth}
\includegraphics[width=\linewidth]{final_bdt_donutExcl_01_mcpc.pdf}
\caption{donutExcl\_01\_mcpc.}
\end{subfigure}%
\begin{subfigure}[t]{0.24\textwidth}
\includegraphics[width=\linewidth]{final_bdt_donutExcl_02_mcpc.pdf}
\caption{donutExcl\_02\_mcpc.}
\end{subfigure}
\caption{Per-variant final cross section with MC-purity correction enabled
(\texttt{mc\_purity\_correction=1}). Nominal (no mcpc) panel included as
reference. The \texttt{donutExcl\_005\_mcpc} panel no longer shows the
negative-yield pathology at high $p_T$.}
\end{figure}

\section*{Full cone, no inner veto ring}

Two additional cross-check variants use the same tower-based isolation
(R=0.4 outer, 120\,MeV threshold, all-calo or ownership-excluded EMCal)
but \emph{skip the inner-$R$ subtraction}. Implemented by setting
\texttt{iso\_emcalinnerr=0.0} in the \texttt{use\_topo\_iso=4/5} block,
so the isolation value is the raw full cone $E_T^\mathrm{iso} =
\texttt{cluster\_iso\_04}$ (or \texttt{\_excl\_04}) with no annulus subtraction.
80\,\% iso cuts rederived via the two-pass workflow:

\begin{table}[H]
\centering
\caption{Derived 80\,\% iso cuts for the full-cone (no-inner-ring) variants.}
\begin{tabular}{lrr}
\toprule
Variant & $b$ [GeV] & $s$ \\
\midrule
donutFull\_noinner & $0.0712$ & $0.01655$ \\
donutExcl\_noinner & $0.2069$ & $0.01889$ \\
\bottomrule
\end{tabular}
\end{table}

\begin{table}[H]
\centering
\caption{Truth purity, closure, and cross-section ratio vs.\ nominal for the
full-cone variants.}
\begin{tabular}{lrrrrr}
\toprule
 & \multicolumn{3}{c}{Truth purity} & \multicolumn{2}{c}{Closure / xsec} \\
\cmidrule(lr){2-4}\cmidrule(lr){5-6}
Variant & $y(9)$ & $y(21)$ & $y(34)$ & $\langle|1-r_\mathrm{cls}|\rangle$ & $\langle|1-r_\mathrm{xs}|\rangle$ \\
\midrule
nom (reference)    & $0.581$ & $0.743$ & $0.797$ & $0.185$ & $0$ \\
donutFull\_noinner & $0.653$ & $0.765$ & $0.827$ & $0.197$ & $0.236$ \\
donutExcl\_noinner & $0.632$ & $0.780$ & $0.843$ & $0.153$ & $0.167$ \\
\bottomrule
\end{tabular}
\end{table}

\textbf{Observations.} Full-cone iso without inner-ring subtraction gives
purity \emph{higher} than nominal at all $p_T$ (by $\sim 0.05$--$0.07$ at
low $p_T$), but closure is only marginally better than nominal for
\texttt{donutExcl\_noinner} ($\langle|1-r|\rangle = 0.153$ vs.\ nominal $0.185$)
and slightly \emph{worse} for \texttt{donutFull\_noinner} ($0.197$).
Cross-section ratios at high $p_T$ are closer to unity ($0.78$--$0.89$) than
some of the donut variants, but mean deviation is larger than the best donut
($R_\mathrm{inner}=0.20$) variants.

\begin{figure}[H]
\centering
\begin{subfigure}[t]{0.32\textwidth}
\includegraphics[width=\linewidth]{purity_sim_bdt_donutFull_noinner.pdf}
\caption{donutFull\_noinner purity.}
\end{subfigure}%
\begin{subfigure}[t]{0.32\textwidth}
\includegraphics[width=\linewidth]{purity_sim_bdt_donutExcl_noinner.pdf}
\caption{donutExcl\_noinner purity.}
\end{subfigure}%
\begin{subfigure}[t]{0.32\textwidth}
\includegraphics[width=\linewidth]{final_bdt_donutFull_noinner.pdf}
\caption{donutFull\_noinner final.}
\end{subfigure}

\vspace{0.3em}

\begin{subfigure}[t]{0.32\textwidth}
\includegraphics[width=\linewidth]{final_bdt_donutExcl_noinner.pdf}
\caption{donutExcl\_noinner final.}
\end{subfigure}
\caption{Full-cone (no inner-ring) variants: per-variant MC purity closure
(top row) and final $d^2\sigma / d\eta\, dp_T$ (bottom row).}
\end{figure}

\end{document}
